\section{Introduction}\label{sec:introduction}
Euler's totient function $\varphi(n)$ counts the integers in $\{1,\ldots,n\}$ that are relatively prime to~$n$. We write
\[
 \mathcal V=\{\varphi(n):n\ge1\},\qquad
 V(x)=\#\{v\in\mathcal V:v\le x\}.
\]
Counting distinct values requires more than counting integers with a
prescribed factorization: different integers can have the same totient.
This distinction is decisive in an asymptotic formula for~$V(x)$.

Throughout the paper logarithms are natural, and $\log_j$ denotes
the $j$-fold iterated logarithm.

\paragraph{The counting problem.}
Since $\varphi(p)=p-1$ for primes~$p$, the prime number theorem gives
$V(x)\ge\pi(x+1)\sim x/\log x$.
Pillai proved that the totient values have density zero
\cite{Pillai1929}. Erd\H{o}s used the normal number of prime factors
of $p-1$ to sharpen the upper bound to
$V(x)\ll_\varepsilon x/(\log x)^{1-\varepsilon}$ for every
$\varepsilon>0$ \cite[Section~2]{Erdos1935}.
His later count of distinct values $\varphi(pr)$ with $p,r$ prime and
$pr\le x$ gives the lower bound
$(1+o(1))x\log_2x/\log x$ \cite[pp.~542--543]{Erdos1945}.
Erd\H{o}s and Hall obtained a much larger factor
$\exp\{a(\log_3x)^2\}$ for every $0<a<1/\log16$
\cite[pp.~1--3]{ErdosHall1976}. Pomerance brought the upper bound onto
the same logarithmic scale \cite[Equation~(1.4)]{Pomerance1986}.
Maier and Pomerance then determined the leading constant
\cite[Section~1]{MaierPomerance1988}:
\[
 V(x)=\frac{x}{\log x}
       \exp\{(C_0+o(1))(\log_3x)^2\},\qquad
 C_0=0.8178146\ldots.
\]

In 1998, Ford refined this to an estimate with only a bounded multiplicative
uncertainty \cite[Theorem~1]{FordTotients2013}:
\begin{equation}\label{eq:classical-order}
 V(x)=\frac{x}{\log x}\exp\left\{
 \begin{aligned}
 &C_0(\log_3x-\log_4x)^2+D_0\log_3x\\
 &\hspace{12pt}-(D_0+\tfrac12-2C_0)\log_4x+O(1)
 \end{aligned}\right\},
\end{equation}
where $D_0=2.1769687\ldots$; both constants are specified in the
notation of Section~\ref{sec:statements} below. Ford also described
the typical prime factorizations of preimages
\cite[Theorems~10 and~16]{FordTotients2013}, which supply the
structural starting point of our proof. All numbered references to
his distribution paper use the revised arXiv version of 14 July 2013.

\paragraph{Fixed dilation.}
Erd\H{o}s and Hall asked whether
\[
 \frac{V(cx)}{V(x)}\longrightarrow c
 \qquad(c>1\text{ fixed})
\]
in the final paragraph of \cite[p.~3]{ErdosHall1976}.
Erd\H{o}s repeated the question in \cite[p.~80]{Erdos1979}.
Ford proved that $V(cx)-V(x)\asymp_c V(x)$ for every fixed $c>1$
\cite[Theorem~4]{FordTotients2013}. This controls the number of
values in a fixed multiplicative interval, but neither it nor
\eqref{eq:classical-order} determines the limiting ratio.

We obtain an explicit asymptotic equivalent and prove the fixed-scale limit, thereby answering this question of Erd\H{o}s and Hall positively. The coefficient in the equivalent retains finite arithmetic information. It is defined in Section~\ref{sec:statements} by bounded prime and integer sums, inclusion--exclusion, and a uniformly convergent limit; its definition does not use~$V$.

\paragraph{The source of the coefficient.}
Order the prime factors of a typical preimage from largest to smallest, and separate a long prefix from a shorter tail. Ford's normal-structure estimates put the double logarithms of the large primes in a simplex. Most of that simplex can be counted by volume. Near its last coordinates the primes remain discrete, however, and several tail factorizations may give the same totient. We retain these tails exactly and take the volume of the union of the regions that they permit. Inclusion--exclusion of this finite union produces the arithmetic coefficient.

Erd\H{o}s's semiprime argument already controls collisions between
products of shifted primes \cite[Lemma~1]{Erdos1945}. The quantitative
shifted-prime normality and layered comparison used here were developed
by Maier and Pomerance \cite[Sections~2--3]{MaierPomerance1988}
and Ford \cite[Section~5]{FordTotients2013}. Here the comparison is
proved under the hypotheses needed for the long prefix, with explicit
costs for recovering the tuple data inside the residual factors.
The resulting control of all ordered prefix collisions, combined with
the exact finite-tail union, converts the volume into a count of
distinct values. Counting the largest prime with common endpoint
parameters first gives regular variation. Identifying the resulting
mass with the arithmetic coefficient then gives the explicit equivalent.
The errors are uniform in the phase; neither conclusion assumes
continuity of the coefficient.

The prefix uniqueness used here is not uniqueness of the entire preimage:
different arithmetic tails may still give the same value. A separate
result on extreme complete fibers gives, for every fixed $\varepsilon>0$,
infinitely many positive integers $v$ with
$\#\{n\ge1:\varphi(n)=v\}>v^{1-\varepsilon}$
\cite[Theorem~1.1]{OpenAITotientFibers2026}. That result concerns
large individual fibers, rather than the asymptotic number of distinct
values, and is not an input to the present proof.

\paragraph{Least preimages.}
For $v\in\mathcal V$, let
\[
 \ell(v)=\min\{n\ge1:\varphi(n)=v\}.
\]
For each fixed positive integer~$k$, the same passage of Erd\H{o}s
\cite[p.~80]{Erdos1979} asks about the number of values $v\le x$ for
which $kx<\ell(v)\le(k+1)x$. Section~\ref{sec:companion} gives a
weighted asymptotic and proves that this count has order $V(x)$
precisely when some totient~$d$ satisfies $\ell(d)>kd$; otherwise
the count vanishes identically. The positive alternative holds for
$k=1,2$. We do not determine which other integers~$k$ satisfy this
seed condition. In particular, the separate question whether
$\ell(d)/d$ is unbounded, posed in \cite[Section~I.9]{Erdos1995},
remains unresolved here.
